\title{Parameter-Efficient Orthogonal Finetuning via Factorization}
Model finetuning fundamentally aims to maintain a high degree of similarity between the pretrained model and its finetuned counterpart, as quantified by specific metrics, so as to retain the knowledge acquired during pretraining. BOFT generalizes the OFT approach and serves as a comprehensive orthogonal finetuning framework. For brevity, we refer to BOFT utilizing $\bm{R}(m, \frac{b}{2})$ as BOFT($m$,$b$), where $b \geq 2$. Drawing inspiration from the Cooley-Tukey fast Fourier transform, which achieves information transmission in an efficient manner, we introduce an orthogonal parameterization that utilizes butterfly architectures. In Figure~\ref{fig:para_eff}, results are averaged across 10 different random seeds. The y-axis measures the error in approximation, while the x-axis quantifies the number of effective, trainable parameters. Importantly, orthogonality imposes specific structural constraints on how nodes can be connected between successive layers. Any orthogonal matrix of dimension $d$ can be expressed by composing butterfly matrices of the form $\bm{B}_{d-1,1}(d)\bm{B}^\top_{d-1,2}(d)\cdots\bm{B}_{1,1}(d)\bm{B}^\top_{1,2}(d)$, where all $\bm{B}_{i,j}(d)$ are butterfly matrices. All together, this framework yields the following forward computation in BOFT:

\begin{equation}
    \begin{aligned}\nonumber
    \bm{z}=\big(\bm{R}(m,b)\cdot\bm{W}^0\big)^\top\bm{x},~~~~\text{s.t.}~\bigg{\{}\bm{R}(m,b)=\prod_{i=1}^m \tilde{\bm{B}}^b_{(d,i)}~~\&~~\underbrace{\big(\tilde{\bm{B}}^b_{(d,j)}\big)^\top\tilde{\bm{B}}^b_{(d,j)}=\tilde{\bm{B}}^b_{(d,j)}\big(\tilde{\bm{B}}^b_{(d,j)}\big)^\top\!=\bm{I}_d}_{\forall j\in [1, m]}\bigg{\}},
    \end{aligned}
\end{equation}

Here, for clarity, we let $\tilde{\bm{B}}^b(d,2^{m - i+1})$ be denoted as $\tilde{\bm{B}}^b_{(d,i)}$, and $\bm{I}_d$ stands for the $d \times d$ identity matrix. Since our orthogonal matrix is represented as a sequence of sparse matrix products, each training step involves repeated matrix multiplications. This leads us to an important inquiry: \emph{Is it possible to devise a dense orthogonal matrix that maintains parameter-efficiency?}

To explore this issue, we introduce a strategy of constructing a dense orthogonal matrix by composing it from a sequence of sparse orthogonal matrices. However, it is important to acknowledge that higher expressivity does not inevitably yield improved performance after finetuning; in practice, fully finetuned models, which can realize arbitrary transformations, may perform suboptimally. BOFT addresses the parameter inefficiency in OFT by leveraging butterfly factorization. Notably, when $m=1$, BOFT($1$,$b$) is equivalent to the block-diagonal OFT~\cite{qiu2023controlling} with block size $b$. The radix-2 Cooley-Tukey procedure systematically decomposes the $N$-point discrete Fourier transform into smaller $\frac{N}{2}$-point transforms, resulting in a pattern known as the butterfly structure. A central challenge in this context is constructing a dense orthogonal matrix without incurring significant parameter overhead. Conventional finetuning strategies typically employ a reduced learning rate, a heuristic technique that implicitly keeps the deviation between the original and finetuned models minimal. BOFT consistently demonstrates superior parameter efficiency when compared to OFT. For example, within a $2\times2$ orthogonal matrix, enforcing a zero in position $(1,1)$ consequently forces a zero at $(2,2)$ as well—that is, omitting one edge might necessitate omitting another. The \emph{butterfly component} is given by $\thickmuskip=2mu \medmuskip=2mu \tilde{\bm{B}}(d,k)=\text{diag}(\bm{B}^F_1(k),\cdots,\bm{B}^F_{d/k}(k))$, a block-diagonal matrix of dimensions $d\times d$ where each block has size $k$, and each $\bm{B}^F_i(k)$ is a butterfly factor as defined in Equation~\ref{eq:but_factor}. A prominent example is LoRA~\cite{hulora2022}, which refines pretrained weights by incorporating the product of two low-rank matrices, delivering strong results on tasks in natural language processing. Specifically, we deploy BOFT in a suite of adaptation scenarios spanning computer vision and natural language processing. The butterfly graph structure intrinsic to the Cooley-Tukey method naturally inspires a technique known as butterfly factorization for decomposing matrices. This rapid advancement has been accompanied by exponential growth in model sizes

(\emph{e.g.}, GPT-3 contains approximately 175 billion parameters). While OFT has exhibited strong generalization capabilities and effective convergence in the context of finetuning text-to-image diffusion models~\cite{qiu2023controlling}, the necessity of employing high-dimensional orthogonal matrices can lead to a substantial number of parameters. As a result, the conventional OFT must set the block size $b=d$ to obtain a dense orthogonal matrix, whereas BOFT achieves this flexibility for arbitrary $b$. BOFT not only encompasses OFT as a particular instance but also generalizes the framework for orthogonal finetuning. This can be formalized via singular value decomposition, where $\bm{W}^0=\bm{U}\bm{\Sigma}\bm{V}^\top$, with $\bm{U}$ and $\bm{V}$ as orthogonal matrices and $\bm{\Sigma}$ as the diagonal matrix of singular values. Building on insights from hyperspherical energy's role in generalization~\cite{liu2018learning,liu2021orthogonal}, \cite{qiu2023controlling} introduces orthogonal finetuning as a compelling alternative, especially for adapting text-to-image diffusion models. The balance between expressivity and regularization is central for robust model finetuning. Various PEFT strategies have been proposed~\cite{houlsby2019parameter,aghajanyan2020intrinsic,hulora2022,edalati2022krona,wang2022adamix,gheini2021cross,zaken2022bitfit,guo2020parameter,sung2021training,ansell2022composable,lester2021power,li2021prefix,vu2022spot,he2021towards,mao2021unipelt,karimi2021compacter,liu2022few,sung2022lst,chen2022parameter,jia2022visual,chen2022adaptformer,zhang2022neural,jie2023fact,lian2022scaling,luo2023towards}, among which reparameterization-based techniques~\cite{aghajanyan2020intrinsic,hulora2022,edalati2022krona} are particularly relevant to this study. \begin{theorem}[Expressivity of BOFT]\label{thm:exp_boft}
    With a fixed block size, BOFT possesses strictly greater expressivity than OFT. For the composite product $\bm{B}_5\bm{B}_4\bm{B}_3\bm{B}_2\bm{B}_1$ to yield a dense $\bm{R}$, every initial node must be able to reach all final nodes in the $6$-th layer. This approach allows us to investigate the parameter-efficiency of OFT from a graphical viewpoint, facilitating the design of efficient matrix decompositions. The proposed information transmission framework draws parallels to concepts from computer networking, particularly studies of network topology efficiency, suggesting that alternative network structures might enable even more parameter-efficient dense orthogonal matrices. From the information transmission standpoint, two main criteria guide our approach: \emph{(i)} \textbf{dense connectivity}, meaning that every node from the input layer is linked, through some sequence of connections, to all nodes in the output layer; and \emph{(ii)} \textbf{minimal free edges}, requiring that the total edge count remains as low as feasible while satisfying the orthogonality requirement. To further improve parameter efficiency within this framework, we draw from the butterfly networks found in the Cooley-Tukey fast Fourier transform algorithm~\cite{cooley1965algorithm}, where nodes can communicate rapidly and efficiently~\cite{klasing1994broadcasting}. Ensuring orthogonality of the block butterfly component $\tilde{\bm{B}}^b(d,2)$ requires parameterizing each $2b\times 2b$ block to be orthogonal. For example, LoRA employs zero-initialization for the additional low-rank matrices. \item We are the first to extend orthogonal finetuning across a diverse array of tasks beyond controllable text-to-image generation~\cite{qiu2023controlling}, thereby demonstrating its promise as a versatile approach to model finetuning. However, this block-diagonal construction imposes a restriction: it essentially groups the dimensions of each neuron (i.e., the columns of $\bm{W}^0$) into $r$ partitions, with each group independently transformed by its corresponding orthogonal block. Furthermore, choosing BOFT($1$,$d$) collapses the formulation into the standard OFT~\cite{qiu2023controlling} using a full, unconstrained orthogonal matrix. For zero initialization, $\bm{R}$ is initialized as the identity matrix. \end{theorem}

This result, stated in Theorem~\ref{thm:exp_boft}, naturally extends BOFT: the resulting orthogonal matrix is expressed as $\thickmuskip=2mu \medmuskip=2mu \bm{R}^G(m_1,b_1,m_2,b_2,l)=\bm{R}_{l,1}(m_1,b_1)\bm{R}_{l,2}^T(m_2,b_2)\cdots\bm{R}_{1,1}(m_1,b_1)\bm{R}_{1,2}^T(m_2,b_2)$, with $\bm{R}_{i,j}^T(m,b)$ denoting an orthogonal matrix of the BOFT type. These connections are mandatory for preserving orthogonality and should be excluded from the edge count, since they correspond to non-trainable elements (in a $d\times d$ orthogonal matrix with $d$ nonzero elements, those values are constrained to $\pm1$). When considering BOFT($\log \frac{2d}{b}$,$b$), the block butterfly matrix $\bm{B}^{\frac{b}{2}}(d)$ serves as $\bm{R}$, thus producing a dense orthogonal matrix $\bm{R}$. As a result, we can now extend orthogonal finetuning's applicability to more general adaptation problems. More generally, for $\bm{R}\in\mathbb{R}^{d\times d}$ to be dense, the collection of factor matrices $\bm{B}_m,\cdots,\bm{B}_1$ must map onto a collection of directed edges across the $d\times (m+1)$ grid, ensuring that each directed edge only links nodes between consecutive layers (that is, adjacent columns), and enabling the transmission of information from each source node at level one to every target node at level $m+1$. For $k>2$, the nonzero structures of $\tilde{\bm{B}}^b(d,k)$ correspond to block-wise permutations of those found in $\tilde{\bm{B}}^b(d,2)$ and thus can be parameterized as orthogonal matrices in the same way. Furthermore, determining an optimal sparsity pattern remains an open problem. In the case where the butterfly parameterization is selected (\emph{i.e.}, $m=\log d, b=2$), the parameter complexity of BOFT is $\mathcal{O}(d\log d)$. Unlike LoRA—which modifies the weights through the addition of a learnable low-rank matrix—OFT employs a multiplicative approach via an orthogonal matrix. \textbf{Inductive bias and generalization in BOFT}. \end{abstract}

\section{Introduction}

Emerging models such as ChatGPT~\cite{brown2020language,chen2021evaluating} and Stable Diffusion~\cite{rombach2022high} have showcased extraordinary generalization capabilities, setting new benchmarks for large foundation models. In order to adapt a pretrained weight matrix, OFT introduces a reparameterization strategy in which the updated weight matrix is expressed as the product of a trainable orthogonal matrix and the fixed pretrained weight matrix. These approaches generally focus on expressing weight matrices directly in terms of the butterfly structure and subsequently training neural networks from initialization. The information transmission perspective offers valuable insights into the sparse structure present in factorization matrices within OFT. Additionally, the matrix factorization inherent in BOFT facilitates a natural form of weight interpolation (refer to Figure~\ref{fig:boft_interpolation}). \textbf{Identity initialization for BOFT}. Notably, after finetuning, the orthogonal matrices optimized via BOFT can be integrated into the pretrained model through multiplication, eliminating additional inference delays. This observation has inspired the development of the Orthogonal Finetuning (OFT) framework~\cite{qiu2023controlling}, in which all neurons within a layer undergo a transformation by a shared orthogonal matrix, thereby ensuring that pairwise neuron angles are exactly preserved during finetuning. Work in~\cite{dao2019learning,dao2019kaleidoscope} demonstrates learning fast linear transforms using butterfly-based parameterizations, while~\cite{chen2022pixelated,dao2022monarch} apply these structures to accelerate neural network training. The orthogonal transformation $\bm{R}(m,b)\in\mathbb{R}^{d\times d}$ is constructed by sequentially multiplying several orthogonal butterfly factors. Traditional Dropout~\cite{srivastava2014dropout} functions effectively with LoRA, but this is not straightforward in BOFT due to the multiplicative structure of its update. A number of prior studies~\cite{dao2019learning,dao2019kaleidoscope,chen2022pixelated,dao2022monarch} have investigated sparse training methods employing butterfly parameterizations. Inspired by the notable structure employed in the Cooley-Tukey algorithm, referred to as butterfly graphs~\cite{cooley1965algorithm}, we note that these graphs can realize dense connectivity between $d$ source and $d$ target nodes with only $\mathcal{O}(d\log d)$ edges. Employing butterfly factorization, a $d$-dimensional dense matrix can be synthesized via the product of $\mathcal{O}(\log d)$ sparse matrices, each with only $\mathcal{O}(d)$ nonzero elements. \textbf{Comparison to butterfly-based sparse training}. To prevent overfitting, LoRA~\cite{hulora2022} includes a Dropout layer that acts on the low-rank update matrices. Consider the case of $\tilde{\bm{B}}(d,2)$, a block-diagonal structure with $2\times 2$ blocks. \section{Orthogonal Parameterization by Butterfly Factorization}

The butterfly arrangement, originally utilized in the Cooley-Tukey algorithm for efficient computation of the fast Fourier transform, is notable for the property that a local adjustment in the frequency domain produces a global effect in the spatial domain—akin to the transmission of information in our context, where any node at the first level may pass data to all nodes in the final level. LoRA imposes a fixed rank across all introduced matrices; to improve parameter allocation across model layers, approaches such as~\cite{zhang2022adaptive,valipour2022dylora,zhang2023increlora} adaptively modulate the rank per layer. Similar correspondences exist for shorter factorizations such as $\bm{B}_3\bm{B}_2\bm{B}_1$ and $\bm{B}_2\bm{B}_1$, aligning the induced graph with their respective sparsity patterns. LoRA attains parameter-efficiency by leveraging low-rank structures, whereas the original OFT model utilizes a block-diagonal configuration for the orthogonal matrix~\cite{qiu2023controlling}; decreasing the block size leads to improved parameter-efficiency for OFT. This distinction is illustrated in Figure~\ref{fig:comp_lora}. Information is first routed according to $\bm{B}_1$, and finally transmitted following $\bm{B}_m$. For instance, to ensure each matrix $\bm{B}_i$ is full-rank—necessary for orthogonality—a set of $d$ edges must connect nodes between adjacent levels to establish a bijective mapping; consequently, at least $d$ edges are present between any two consecutive layers (see the scenario depicted in Figure~\ref{fig:info_trans} for $d=4$). For instance, the representational capacity of BOFT($9$,$2$) surpasses that of BOFT($1$,$16$) while requiring significantly fewer parameters. The butterfly structure’s prevalence in efficient linear operations, such as discrete Fourier and cosine transforms, supports our hypothesis that this structured parameterization induces an implicit inductive bias conducive to improved generalization and mitigated overfitting. In order to maintain the orthogonality of $\bm{R}$ during fine-tuning, we utilize the Cayley parameterization as described in~\cite{qiu2023controlling,liu2021orthogonal}, namely $\bm{R}=(\bm{I}+\bm{Q})(\bm{I}-\bm{Q})^{-1}$, where $\bm{Q}$ is a skew-symmetric matrix ($\bm{Q}=-\bm{Q}^\top$). The depicted graph is constructed by sequentially unfolding the matrix multiplication. Specifically, an entry $R_{ij}$ of $\bm{R}$ being zero signals that there is no route for information from node $j$ to node $i$, while a nonzero value permits such flow. Orthogonality of $\tilde{\bm{B}}(d,2)$ can be assured by applying a Cayley transform, or alternatively, by expressing each block as a $2$-dimensional rotation, in accordance with \cite{liu2021orthogonal,qiu2023controlling}. This formulation recasts the design of dense, parameter-efficient orthogonal matrices as an information routing problem on a graph with a grid topology. Compared to LoRA, orthogonal finetuning generally achieves superior preservation of spectral characteristics because it exactly maintains the spectral norm of the pretrained parameter matrix $\bm{W}^0$. To facilitate practical matrix factorization discovery, we introduce a graphical information transmission framework. & \bm{0} \\
    \bm{0}\!\!\!\!\! Beyond these, butterfly matrices also contribute to efficient matrix-vector multiplication~\cite{michielssen1996multilevel,o2010algorithm}, facilitate data-sparse matrix representations~\cite{li2015butterfly}, and support network information transmission~\cite{klasing1994broadcasting,li2003linear,rajasingh2016transmission}. Analyzing the resulting nonzero structure in a composite orthogonal matrix highlights the strength of the information transmission approach for OFT, because it allows us to focus solely on the pattern and position of nonzero, trainable entries in $\bm{R}$, independent of their actual numerical values. Our reasoning is based on the butterfly structure’s capacity to represent a wide array of classical linear transforms—a capability not shared by OFT’s block-diagonal construction. Our experiments confirm that BOFT delivers robust finetuning performance on a range of large-scale models, including language, vision, and text-to-image generation, further supporting the utility of BOFT as a broadly applicable finetuning technique. As a result, building these large-scale models from the ground up has become exceedingly difficult and costly. Addressing the training runtime efficiency of BOFT is a subject for future research. \section{Intriguing Insights and Discussions}

\textbf{Expressivity of BOFT}. \textbf{Butterfly structures}. Although OFT shows promising generalization, its reliance on orthogonal matrices entails a considerable number of learnable parameters, limiting parameter efficiency. This information transmission view offers a flexible design space for sparse matrix factorizations, allowing us to keep the parameter count low while maintaining sufficient expressive power to form dense matrices. \input{rebuttal-results}

\section{Concluding Remarks and Limitations}

This work presents Orthogonal Butterfly, a universal parameter-efficient finetuning approach for foundation models leveraging the butterfly matrix structure. Concretely, let us consider the generation of a dense matrix $\bm{R}\in\mathbb{R}^{d\times d}$ expressed as the product of $m$ square matrices: $\thickmuskip=2mu \medmuskip=2mu \bm{R}=\bm{B}_m\bm{B}_{m-1}\cdots\bm{B}_1$. Such matrices have found applications in orthogonal matrix parameterization, mitigating the need for pivoting in Gaussian elimination and enhancing computational efficiency~\cite{parker1995random}; improving the stability of recurrent neural network training~\cite{jing2017tunable}; and kernel approximation~\cite{munkhoeva2018quadrature}. To mitigate this, OFT utilizes a block-diagonal matrix structure to limit parameter count. \textbf{Orthogonal finetuning as learning bilinear similarity}. Nevertheless, BOFT has certain limitations. The BOFT approach is equivalent to learning a bilinear similarity function, where the similarity score between the $i$-th neuron (that is, the $i$-th column of $\bm{W}^0$), $\bm{w}^0_i$, and input $\bm{x}$ is computed as $\bm{w}^0_i\bm{R}\bm{x}$. For instance, the network shown in Figure~\ref{fig:info_trans} achieves dense interconnections using $10$ links, whereas the butterfly topology requires only $8$. With the advent of increasingly sophisticated foundation models (\emph{e.g.}, \cite{brown2020language,radford2021learning,rombach2022high,kirillov2023segment}), there has been growing interest in finetuning these models for downstream applications while minimizing the number of additional parameters~\cite{treviso2023efficient,ding2023parameter,lialin2023scaling}. Furthermore, its effectiveness across a broader set of adaptation scenarios—beyond those explored in~\cite{qiu2023controlling}—remains to be established. The tractability and effectiveness of low-rank parameterization have led to widespread adoption~\cite{dettmers2023qlora,zi2023delta,chavan2023one}. Because $\bm{R}(m,b)$ in BOFT typically corresponds to a structured subspace of the full orthogonal group—thereby imposing additional restrictions on the space of candidate hypotheses—BOFT inherently introduces a form of inductive bias. Consequently, it has become crucial to develop strategies for adapting these extensive models to new tasks in an efficient manner. Among these, model finetuning stands out for its simplicity, effectiveness, and the advantage of not increasing inference latency. For BOFT, all butterfly components are initially set to identity matrices (i.e., setting the associated skew-symmetric matrices to zero prior to applying the Cayley transform). &\bm{B}_2(\frac{d}{2})
    \end{bmatrix}= \tilde{\bm{B}}(d,d)\tilde{\bm{B}}(d,\frac{d}{2})\cdots\tilde{\bm{B}}(d,2),
\end{aligned}
\end{equation}

Here, $\bm{B}_1(\frac{d}{2})$ and $\bm{B}_2(\frac{d}{2})$ represent two butterfly matrices each of dimension $\frac{d}{2}$. Utilizing the butterfly matrix to parameterize an orthogonal matrix is now straightforward: it is sufficient to guarantee that all the multiplicative components $\tilde{\bm{B}}(d,k)$ of the butterfly matrix $\bm{B}(d)$ are orthogonal for all $k$. Next, when $d=2^N$, we recursively construct the $d$-dimensional \emph{butterfly matrix} $\bm{B}(d)\in\mathbb{R}^{d\times d}$ as

\begin{equation}
\begin{aligned}
    \!\!\bm{B}(d)=\tilde{\bm{B}}(d,d)\!\cdot\!\begin{bmatrix}
    \bm{B}_1(\frac{d}{2})\!\!\!\!\! \end{itemize}

\section{Related Work}

\textbf{Parameter-efficient finetuning (PEFT)}. Each color-specific curve corresponds to a BOFT configuration with a fixed block size; the leftmost data point on each curve reflects the performance of BOFT($1$,$b$) (\emph{i.e.}, the standard block-diagonal OFT using block size $b$). Formally, OFT seeks to learn an orthogonal matrix $\bm{R}\in\mathbb{R}^{d\times d}$ for a given pretrained linear layer $\bm{W}^0\in\mathbb{R}^{d\times n}$, adjusting the computation in the forward pass from $\bm{z}=(\bm{W}^0)^\top\bm{x}$ to $\bm{z}=(\bm{R}\bm{W}^0)^\top\bm{x}$, with $\bm{x}\in\mathbb{R}^{d}$ as the input and $\bm{z}\in\mathbb{R}^{n}$ as the output vector. Nonetheless, OFT typically incurs a higher parameter count than LoRA, prompting efforts to enhance its efficiency. Central to our method is the insight that greater parameter efficiency can be achieved by constructing a dense orthogonal matrix through the multiplication of several sparse orthogonal matrices. Thus, the set of feasible edge patterns is sometimes expanded or limited by the orthogonality constraints. Generally, there are three prevalent strategies for efficient task adaptation: \emph{model finetuning}~\cite{hulora2022,zaken2022bitfit,guo2020parameter,raffel2020exploring,qiu2023controlling,zhang2022adaptive,chavan2023one}, wherein the architecture remains intact but a selection of the original parameters are updated; \emph{adapter tuning}~\cite{rebuffi2017learning,houlsby2019parameter,pfeiffer2020adapterfusion,lin2020exploring,he2021towards}, which involves introducing additional trainable components to the model; and \emph{prompt tuning}~\cite{li2021prefix,lester2021power}, where trainable prefixes are prepended to the inputs. Through this lens, the butterfly structure emerges as an effective strategy for sparse orthogonal matrix decomposition. Efficient adaptation of these models to downstream applications is therefore a core concern. This is in contrast to the original OFT with a full orthogonal parameterization, which necessitates $\mathcal{O}(d^2)$ parameters, or the block-diagonal OFT with $r$ blocks, which involves $\mathcal{O}(bd)$. By enriching the finetuning parameter space with structural constraints, BOFT facilitates the pursuit of an optimal balance between expressive power and inductive bias. For continued progress across the field, it is vital to develop methods that can effectively adapt pre-trained foundation models to specific tasks without necessitating full retraining. This butterfly configuration is also prominent as a network topology in computing systems~\cite{solihin2015fundamentals,li2011linear} for its efficient communication capabilities. A comparison with LoRA~\cite{hulora2022} is instructive: both low-rank and sparse matrix approaches represent prominent categories of structured matrices~\cite{candes2011robust}, each offering distinct avenues for parameter-efficient learning. Notably, both block-diagonal and butterfly structures incorporate sparsity to orthogonal matrices, thereby lowering the number of required learnable parameters. The primary contributions of our work are summarized below:

\begin{itemize}[leftmargin=*,nosep]
    \item We address parameter-efficient orthogonal finetuning by introducing a novel framework grounded in information transmission. The arrangement of nonzero entries in any butterfly component can be interpreted as a rearrangement of those in $\tilde{\bm{B}}(d,2)$; as a result, all such components can be analogously parameterized as orthogonal matrices. \textbf{Spectral properties}. \textbf{Multiplicative Dropout for BOFT}. It is also unclear whether the butterfly topology is the most optimal for information transmission. This structural choice allows us to target a more compact set of hypotheses within the orthogonal group for downstream applications. In contrast to the block-diagonal approach adopted by OFT, which balances model expressivity and regularity, BOFT employs a butterfly structure to provide a more gradual transition between full orthogonal matrices (promoting expressivity) and identity matrices (encouraging regularity). To conceptualize this factorization, we recast dense orthogonal matrix generation as an information transmission process. Consequently, decomposing $\bm{R}$ into multiple factors $\bm{B}_m\bm{B}_{m-1}\cdots\bm{B}_1$ can be viewed as orchestrating successive information exchanges along a sequence of graphs induced by the $\bm{B}_i$ matrices. Recognizing the inherent structure of weight matrices, a more theoretically motivated strategy seeks to conserve the relational structure within these matrices—specifically, the angles formed between pairs of neurons. This process can be algebraically represented by a product of sparse matrices, collectively called a butterfly matrix. Nonetheless, this approach imposes its own limitations: the orthogonal matrices possess fixed sparsity patterns, and each block operates its transformation independently. This approach results in a computationally efficient orthogonal parameterization that constructs the larger matrix from numerous $2\times 2$ orthogonal components. As illustrated in Figure~\ref{fig:blk_diag}, the original OFT exhibits a block-diagonal form for $\bm{R}$. \section{An Information Transmission View on Orthogonal Finetuning}

We begin by outlining key concepts of OFT. Although this technique is effective in practice, dividing a neuron's dimensions by index alone lacks a principled basis, highlighting an inherent limitation and thereby motivating interest in dense orthogonal matrices. Such multi-level compositions are often described as the kaleidoscope hierarchy~\cite{dao2019kaleidoscope}. In the following, we discuss how the butterfly configuration enhances parameter efficiency. Generally, the BOFT($m$,$b$) structure provides $\frac{1}{2}(b-1)dm$ adaptable parameters for tuning a linear mapping of dimensions $d \times n$. A straightforward fully-connected scheme between two layers involves $\mathcal{O}(d^2)$ connections (that is, represented by a dense orthogonal matrix), resulting in $d^2-d$ connections (for $d=4$, this amounts to 12 connections). To overcome this challenge, we introduce a multiplicative Dropout scheme for BOFT. In contrast, \cite{dao2022monarch} examines a finetuning approach in which pretrained weights are first projected onto modified butterfly matrices, with subsequent optimization of these projected elements for specific tasks. To overcome this limitation, we analyze OFT from an information transfer standpoint, identifying several essential criteria that promote greater parameter-efficiency. In general, choosing a smaller block size $b$ and larger $m$ for BOFT results in enhanced parameter efficiency; e.g., BOFT($6$,$4$) achieves an approximation error comparable to that of BOFT($2$,$16$), but utilizes substantially fewer parameters. Although the block-diagonal sparsity enhances parameter efficiency, it may inadvertently introduce undesirable inductive biases (such as diminished expressivity or the inability to approximate canonical linear transformations). Integrating this parameterization with OFT yields a new, highly parameter-efficient finetuning approach we term Orthogonal Butterfly~(BOFT). By ensuring orthogonality in each sparse component, this results in an orthogonal parameterization requiring merely $\mathcal{O}(d \log d)$ parameters—a substantial reduction compared to conventional approaches. In this framework, the optimization seeks a bilinear similarity matrix $\bm{R}$ that is subject to the strong regularization constraint of orthogonality, closely linking BOFT with concepts from distance metric learning~\cite{xing2002distance} and the general bilinear form~\cite{roman2005advanced}. Our approach circumvents this limitation by assembling a dense orthogonal matrix through the product of several factorized sparse matrices. This work examines Orthogonal Finetuning (OFT), a structured finetuning strategy aimed at adapting large models for specific downstream tasks. To further enhance parameter-efficiency, the orthogonal matrix $\bm{R}$ is constructed to be block-diagonal by expressing it as $\text{diag}(\bm{R}_1,\bm{R}_2,\cdots,\bm{R}_r)$, with each $\bm{R}_i\in\mathcal{R}^{b\times b}, \forall i$, serving as a small orthogonal matrix such that $br=d$. To investigate this, we perform a numerical experiment that attempts to approximate a randomly sampled dense orthogonal matrix~\cite{anderson1987generation} of dimensions $512\times 512$. This methodology leverages the fundamental idea that computational expense may be substituted for memory efficiency, reducing the number of learnable parameters. Directly parameterizing a dense $d$-dimensional orthogonal matrix would conventionally demand $\mathcal{O}(d^2)$ parameters, rendering it impractical for large $d$. For a cohesive and transparent way to analyze the sparsity structure involved in orthogonal matrix factorizations, we conceptualize the generation of a dense orthogonal matrix within OFT as a process analogous to information propagation. Since $\bm{R}(m,b)$ is constructed by $m$ orthogonal butterfly factors—each of which may be arranged into $2b \times 2b$ block-diagonal orthogonal matrices—the proposed Dropout randomly selects $p_1$ percent of the butterfly factors and, within each selected factor, $p_2$ percent of its diagonal blocks, then replaces these selected blocks by identity matrices. This can be interpreted as an information transfer over a network laid out in a $d\times (m+1)$ grid, as depicted in Figure~\ref{fig:info_trans}. Additionally, the use of sparse matrix factorization within BOFT could impart further implicit inductive biases~\cite{gunasekar2017implicit,li2018algorithmic,liu2019neural}. Let $k\geq 2$ be a power of $2$; we define the \emph{butterfly factor} $\bm{B}^F(k)$ by

\begin{equation}\label{eq:but_factor}
\begin{aligned}
    \bm{B}^F(k)=\begin{bmatrix}
    \text{diag}(\bm{d}_1) & \text{diag}(\bm{d}_2) \\
    \text{diag}(\bm{d}_3) & \text{diag}(\bm{d}_4)
    \end{bmatrix}\in\mathbb{R}^{k\times k},
\end{aligned}
\end{equation}

where for each $i$, the vector $\bm{d}_i\in\mathbb{R}^{\frac{k}{2}}$. Here, each matrix $\bm{B}_i$ describes connections from nodes at the $i$-th layer to those in the $(i+1)$-th layer, with each matrix entry $(j_1,j_2)$ indicating the presence (nonzero) or absence (zero) of a directed edge from the $j_2$-th node at level $i$ to the $j_1$-th node at level $i+1$. Distinct from these, BOFT introduces an alternative finetuning methodology that leverages layer-shared weight matrices to transform weights. Our approach departs from earlier work by specifically utilizing butterfly structures to improve the parameter-efficiency of OFT. However, restricting our attention to the product $\bm{B}_4\bm{B}_3\bm{B}_2\bm{B}_1$, relating nodes from the first to fifth layers, reveals that node 1 cannot transmit to node 3, corresponding to a zero entry at $(3,1)$ in the sparsity structure of $\bm{B}_4\bm{B}_3\bm{B}_2\bm{B}_1$. In both OFT and BOFT, the application of an orthogonal transformation $\bm{R}$ from the left yields updated weights of the form $\bm{R}\bm{U}\bm{\Sigma}\bm{V}^\top$, leaving the maximum singular value, that is, the spectral norm of $\bm{W}^0$, unaltered. From this perspective, constructing a dense orthogonal matrix by multiplying sparse matrices corresponds to disseminating information across a grid-structured graph. This viewpoint is motivated by recognizing that a dense $d\times d$ matrix essentially encodes all-to-all connections from $d$ input nodes to $d$ output nodes. Following \cite{chen2022pixelated}, this construction is further generalized to define a block butterfly component $\tilde{\bm{B}}^b(d,k)$, in which each scalar entry in $\bm{d}_i$ is instead replaced by a $b\times b$ matrix. In transfer learning, it is common practice to initialize with the pretrained model’s parameters so that finetuning does not result in substantial divergence from the original model. Further mathematical insights are discussed in Appendix~\ref{app:sn_property}. While the butterfly framework combined with permutations can exactly reconstruct several classical fast linear transforms~\cite{dao2019learning,dao2019kaleidoscope} (\emph{e.g.}, fast Fourier transform, Hadamard transform), the extent to which an orthogonal butterfly matrix can closely approximate an arbitrary orthogonal matrix has not been fully established. We posit that the particular structural bias arising from butterfly factorization enhances generalization capabilities, as it shares structural motifs with several foundational linear transforms, including the discrete Fourier, sine/cosine, and Hadamard transforms~\cite{dao2019kaleidoscope}. Exploiting this efficient parameterization, we introduce Orthogonal Butterfly~(BOFT), a novel parameter-efficient finetuning technique. Notably, if $m_1=m_2=\log d$, $b_1=b_2=2$, and $l=d-1$, this composition $\bm{R}^G(m_1,b_1,m_2,b_2,l)$ is sufficient to span the entire orthogonal group. The primary reason for utilizing orthogonal transformations in fine-tuning is to maintain the pairwise angular relationships among neurons~\cite{liu2018learning,liu2021orthogonal,qiu2023controlling}, thereby ensuring the pretrained semantic information is largely retained. Since orthogonal matrices generally demand fewer parameters compared to unconstrained full matrices, orthogonality introduces additional interdependencies among edge selections. \item We present several theoretical arguments to elucidate why BOFT can drastically decrease the parameter count needed for learning while preserving both expressive power and stability during training. \item Drawing inspiration from the butterfly arrangement seen in the Cooley-Tukey algorithm, we introduce Orthogonal Butterfly, a new parameter-efficient finetuning scheme for orthogonal matrices within the context of our information transmission framework. This advantage arises because the butterfly matrix constrains the representation to a more structured subset within the orthogonal group (in comparison with the block-diagonal structure), implying that BOFT is provably more expressive than OFT for the same block size. 
\begin{abstract}
As large foundation models become increasingly prevalent, the high costs associated with training them from scratch pose significant challenges. In these settings, BOFT consistently surpasses leading parameter-efficient methods, highlighting its strength in both efficiency and generalizability. The central objective is to build a dense orthogonal matrix by combining $m$ sparse orthogonal matrices, all while minimizing the number of essential trainable parameters. Comprehensive empirical evaluations demonstrate the effectiveness of BOFT in adapting a range of large models—including vision transformers, language models, and text-to-image diffusion systems—to a variety of vision and language downstream tasks. This technique, referred to as OFT, learns an orthogonal transformation within each layer, yielding better generalization and more stable training dynamics compared to LoRA. As an explicit instance in Figure~\ref{fig:info_trans}, we examine the decomposition $\bm{R}=\bm{B}_5\bm{B}_4\bm{B}_3\bm{B}_2\bm{B}_1$, highlighting both its sparsity layout and corresponding graph structure. As illustrated in Figure~\ref{fig:info_trans}, an example of a valid matrix decomposition requires just $10$ connections, which is fewer than required by a fully dense orthogonal matrix. Prior studies have shown that such invariance substantially supports both stability in training and improved generalization~\cite{miyato2018spectral,yoshida2017spectral}. Due to the nature of forming the final orthogonal matrix as a product of several orthogonal matrices, the method introduces a modest increase in training computational cost relative to OFT. Although this structure decreases the parameter count, it does not permit $\bm{R}$ to be a fully dense matrix.